\section{In-Sample Results and Cross-Verification}\label{sec:insample}

\subsection{Estimating the inputs}

$\muhat$ is the column-wise sample mean of the return panel and $\Sig$ its sample covariance
($\text{ddof} = 1$), symmetrized. The square root $\Ssqrt$ is built by \textbf{eigendecomposition
rather than Cholesky factorization}: a sample covariance estimated with $N = 98$ from $T = 965$
observations is positive \emph{semi}definite at best and in floating point carries a few
marginally negative eigenvalues, on which Cholesky --- requiring strict positive definiteness ---
simply fails. The implementation computes $\Sig = V\operatorname{diag}(\ell)V^\top$, clips
negative $\ell_i$ to zero, and sets $\Ssqrt = V\operatorname{diag}(\sqrt{\ell})V^\top$,
reconstructing $\Sig$ to $\norm{\Ssqrt\Ssqrt - \Sig}_F / \norm{\Sig}_F = 1.831 \times 10^{-15}$.

\subsection{Solver behaviour and the effect of $\lambda$}

\Cref{tab:demo} reports two representative solves. In both, $\one^\top w = 1$ holds to
$10^{-12}$ and the reported zeros are \textbf{exact} zeros produced by the proximal step, not
small values below a display threshold. That is the practical payoff of \cref{prop:jointprox}:
the sparsity is real, so the count of nonzero weights is a count of positions the investor would
actually open.

\begin{table}[htbp]
\centering
\caption{Two representative solves on the full return panel.}
\label{tab:demo}
\small
\begin{tabular}{@{}lrrrr@{}}
\toprule
$(\kappa,\gamma,\lambda)$ & Iterations & Objective & Active $(|w_i| > 10^{-4})$ & Exact zeros \\
\midrule
$(1,\, 5,\, 0.01)$ & 528  & 0.015302 & 17 / 98 & 81 / 98 \\
$(0,\, 5,\, 0.10)$ & 2113 & 0.099071 & 10 / 98 & 88 / 98 \\
\bottomrule
\end{tabular}
\end{table}

\begin{figure}[htbp]
\centering
\begin{subfigure}[t]{0.44\textwidth}
  \centering
  \includegraphics[width=\textwidth]{fig2_convergence.png}
  \caption{Objective against iteration count.}
  \label{fig:convergence}
\end{subfigure}
\hfill
\begin{subfigure}[t]{0.44\textwidth}
  \centering
  \includegraphics[width=\textwidth]{fig3_sparsity_path.png}
  \caption{Active positions against $\lambda$.}
  \label{fig:sparsity}
\end{subfigure}
\caption{Left: convergence; the non-monotone early phase is expected of a subgradient method and
is why the solver returns the best iterate. Right: the sparsity path.}
\label{fig:solver}
\end{figure}

\Cref{fig:sparsity} traces active positions as $\lambda$ sweeps a logarithmic grid. The count
falls as $\lambda$ rises but \textbf{not strictly monotonically}, for a structural rather than
numerical reason: $\lambda$ enters only through the threshold $t = \alpha_k\lambda$, and raising
$t$ also shifts the multiplier $\nu^\star$ enforcing the budget constraint in
\cref{prop:jointprox}, so a larger $\lambda$ can move $z_i + \nu^\star$ far enough for a
coordinate to re-enter the active set even as others leave. The trend is unambiguous; the
monotonicity is not, and the figure is reported as it is rather than smoothed.

\subsection{The effect of the robust term}

Raising $\kappa$ from $0$ to $1$ with $\gamma, \lambda$ fixed \textbf{increases} both the number
of active positions and the Herfindahl--Hirschman concentration index: the portfolio becomes
\emph{more} concentrated, where the folk intuition --- robust means cautious means diversified
--- predicts the opposite. The mechanism is the degree-one homogeneity of
\cref{eq:robust-derivation}: because $\kappa\norm{\Ssqrt w}_2$ is a norm rather than a squared
norm, its subgradient $\kappa\,\Sig w / \norm{\Ssqrt w}_2$ does \emph{not} shrink in magnitude as
$\norm{w}$ shrinks, unlike the variance gradient $2\gamma\Sig w$, so it exerts a scale-independent
pull in a fixed direction --- which, given a covariance matrix dominated by a single market
factor, is served more efficiently by concentrating weight along a low-variance combination of a
few names than by spreading it across many correlated ones. This is an empirical observation on
one dataset with a plausible mechanism, \emph{not} a theorem about every $\Sig$.
\Cref{app:robust} gives the figure.

\subsection{Cross-verification against an interior-point solver}

The method of \cref{sec:method} is first-order, hand-written, and specific to this problem.
CLARABEL \citep{clarabel2024}, invoked through CVXPY \citep{diamond2016}, is a general-purpose
second-order interior-point solver operating on the conic reformulation \cref{eq:socp}. They
share no code, no algorithmic family, and no failure modes, and \cref{sec:correctness}
established that the optimal value is unique --- so two correct solvers must agree on it.
Agreement on the objective is the weaker test; agreement on the \textbf{active set} is far
stronger, because the support is a discrete object that a numerically close but structurally
wrong solution --- such as one from the scheme of \cref{sec:proxbug} --- will generally get wrong
even while landing near the right objective value. That is exactly how the bug was caught.

\begin{table}[htbp]
\centering
\caption{Hand-written solver against CVXPY/CLARABEL. Jaccard $= 1.00$ means the supports are
identical. The full table, with $\ell_\infty$ weight distances and iteration counts, is in
\cref{app:verification}.}
\label{tab:cvxpy}
\small
\begin{tabular}{@{}rrrrrrrr@{}}
\toprule
$\kappa$ & $\gamma$ & $\lambda$ & $f_{\text{hand}}$ & Rel.\ gap
  & Active (hand) & Active (CVXPY) & Jaccard \\
\midrule
1   & 5  & 0.01  & 0.0153016 & $1.92 \times 10^{-7}$  & 17 & 17 & 1.00 \\
0   & 5  & 0.10  & 0.0990708 & $7.75 \times 10^{-6}$  & 10 & 10 & 1.00 \\
1   & 5  & 0.001 & 0.0060877 & $2.44 \times 10^{-7}$  & 39 & 39 & 1.00 \\
2   & 10 & 0.05  & 0.0607131 & $-1.66 \times 10^{-8}$ & 16 & 16 & 1.00 \\
0.5 & 1  & 0.02  & 0.0225015 & $1.82 \times 10^{-6}$  & 18 & 18 & 1.00 \\
1   & 5  & 0     & 0.0045395 & $5.14 \times 10^{-7}$  & 98 & 98 & 1.00 \\
\bottomrule
\end{tabular}
\end{table}

The result is unambiguous: the \textbf{maximum relative objective gap is $0.000775\%$} and the
\textbf{active sets match exactly on all six of six parameter sets}, every run reporting
convergence; \cref{fig:proxcvx} in \cref{app:verification} shows the individual weights lying on
the line $y = x$.

The last row is a deliberate control: with $\lambda = 0$ the $\ell_1$ penalty is off and both
solvers return a fully dense portfolio of 98 active positions. Without it, identical supports
could in principle reflect nothing more than two methods trivially zeroing the same coordinates;
the dense case shows the agreement survives when there is no sparsity to agree about.

One detail deserves disclosure. CVXPY rejects a covariance matrix that is not strictly positive
semidefinite in floating point, so the code wraps $\Sig$ in \texttt{cp.psd\_wrap} to suppress
that check. This relaxes a numerical tolerance --- the same concession as clipping eigenvalues
when forming $\Ssqrt$ --- and does not alter the data passed to either solver.
